\chapter{The Language}
\label{ch:language}

\section{Introduction}

A compiler for many targets has to know what each target does. The usual answer is to write it
down: a driver, a machine description, a cost model per part. Gnascor answers the other way. It
asks the part, records what the part answered and what each answer cost, and derives every table
it needs from those records.

Every host is reduced to one interface, a question with a yes or no answer put to an address. A
processor, a database, a web service and a language model all have addresses, and all of them
spend time, compute or tokens to answer. In the terms of software testing each host is a black
box behind a test oracle, a component that only says whether a check holds \cite{src:Howden-1978}.
The language never models what is inside the box. It reads the answers and what they cost.

This chapter states the language and its principle. Chapter~\ref{ch:protocol} states the query
protocol as definitions and propositions, Chapter~\ref{ch:states} the state machine its keywords
name, Chapter~\ref{ch:measured} the measurements, and Chapter~\ref{ch:hypotheses} the hypotheses
they support and the problems left open.

\section{Two faces}

A program has two faces that compile to the same machine code. The high order face writes plain
statements: \emph{a equals something; b equals something; evaluate a is identical to b}. The
assembly face writes the same program short: \texttt{a=0;b=1; x = a==b}. Both compile to two
reads, one target and one store.

A block declares its face, and a statement mixes the two only inside an explicit escape,
\texttt{\_\_gsm\_\_(...)}. This is inline assembly with a stricter rule: the face is checked, and
nothing crosses between faces unmarked.

\section{The data model}

The language keeps what it learns in files grouped by stem. Files sharing a stem describe one
member, a part or a set, and the suffix says which view of the member a file holds. A member
carries only the views it has answers for. Table~\ref{tab:kfiles} gives each view beside the
field's nearest term.

\begin{table}[h]
\centering
\small
\begin{tabular}{@{}l l l@{}}
\toprule
suffix & holds & field term \\
\midrule
\texttt{.ksc} & the language map & language specification \\
\texttt{.krs} & the coherence rules & rule set; instruction selection table \\
\texttt{.kcr} & information at or near its Kolmogorov complexity & compressed form \\
\texttt{.knf} & a measured noise floor & baseline; noise floor \\
\texttt{.kcs} & what rebuilds information & generator; decoder \\
\texttt{.kdm} & the hardware map & machine description \\
\bottomrule
\end{tabular}
\caption{The views of a member. The suffixes are the language's own.}
\label{tab:kfiles}
\end{table}

A hardware map holds a general block, valid for any part of a class, and under it a block
specific enough to be best on one part and no other. A driver written by hand is the general
block only, since a person writes it once. Here neither block is written: both are read off the
part, and keeping the specific one costs nothing.

One further view holds every question put to a member and every answer, with costs, refusals and
censored samples marked, and the best path per problem read from those same answers. In the
field this is a measurement log. The language has not yet named its suffix.

\section{Information is coherence}

The views carry Kolmogorov's name because the language rests on one principle. A description at
its Kolmogorov complexity has no redundancy: no part of it can be predicted from the rest
\cite{src:Kolmogorov}. A system at coherence has the same property seen from outside. Its parts
agree, none spends effort undoing another, and the friction between them is at its floor. Noise
costs and carries nothing; information aligns and carries.

The length of a shortest description cannot be computed. Friction can be measured, in time,
compute and latency, on any host that answers at all, and the language takes it as its reading of
information content. This is the minimum description length principle read from the side of the
machine \cite{src:Rissanen}: the best model of data describes it most briefly, and the best
arrangement of a computation runs on the host with least friction.
Chapter~\ref{ch:hypotheses} asks what the measured friction tracks, and finds it closer to
Bennett's logical depth \cite{src:Bennett} than to description length.

The crystal, the floor and the construction set are the three sides of the principle. A crystal
is information at its complexity, a floor is the part that carries nothing, and a construction set
rebuilds information from its primitives.

\section{Keywords}

The language's keywords are four-letter mnemonics. They name the states of a branch pair and the
transitions between them, and each maps to a concept the field already has a word for.
Table~\ref{tab:mnemonics} gives the map; Chapter~\ref{ch:states} gives the machine, and
Chapter~\ref{ch:terms} the sources.

\begin{table}[h]
\centering
\small
\begin{tabular}{@{}l l l@{}}
\toprule
keyword & what it names & field term \\
\midrule
dual & both sides answer inside the timeout & active-active \\
lead, rite & one side answers & active-passive \\
void & neither side answers & outage \\
busy & both answer, both slow & saturation \\
wait & one side lags & latency \\
blok & a side refuses & error \\
gray & a side not yet asked & unknown; Belnap's None \\
pass, back & the answering side changes over & failover, failback \\
vent & a saturated pair drops one side & load shedding \\
fuse & both sides lost at once & circuit breaker tripped \\
jamm & saturation turns into refusal & gridlock \\
sprk & from nothing to both & recovery \\
sync, gray & entered and left by marked labels & superstate \\
the clock & the stream of states and labels & health status stream \\
\bottomrule
\end{tabular}
\caption{The keywords beside the field's terms.}
\label{tab:mnemonics}
\end{table}
